% !TEX root = ../../ctfp-print.tex

\lettrine[lhang=0.17]{本}{書の第一部では、}圏論とプログラミングはどちらも合成可能性に関わるものだと論じた。
プログラミングでは、問題を扱えるレベルまで分解し、各部分問題を順番に解いて、
解を下から積み上げて再合成する。大まかに言えば、その方法は二つある。
コンピュータに何をするかを伝える方法と、どうするかを伝える方法だ。
一方は宣言的、もう一方は命令的と呼ばれる。

これは最も基本的なレベルでも見られる。合成自体を宣言的に定義することができる。
つまり、\code{h}は\code{f}の後に\code{g}を合成したものである、というように：

\src{snippet01}
あるいは命令的に定義することもできる。つまり、まず\code{f}を呼び出し、その呼び出し結果を記憶し、
その結果で\code{g}を呼び出す、というように：

\src{snippet02}
プログラムの命令的なバージョンは通常、時間的に順序付けられた一連の動作として記述される。
特に、\code{g}の呼び出しは\code{f}の実行が完了する前には起こり得ない。
少なくとも、それが概念的なイメージだ――\emph{必要呼び}の引数渡しを持つ遅延評価言語では、
実際の実行は異なる順序で進むこともある。

実際、コンパイラの賢さ次第では、宣言的コードと命令的コードがどのように実行されるかに
ほとんど差がない場合もある。しかし二つの方法論は、問題解決へのアプローチの仕方や、
生成されるコードの保守性とテスト可能性において、時には大きく異なる。

主要な問いはこうだ：問題に直面したとき、宣言的アプローチと命令的アプローチのどちらかを
常に選べるのか？そして、宣言的な解法が存在する場合、それは常にコンピュータコードに変換できるのか？
この問いへの答えは明らかではなく、もし答えが見つかれば、宇宙への理解を革命的に変えるだろう。

\begin{wrapfigure}{R}{0pt}
  \includegraphics[width=0.5\textwidth]{images/asteroids.png}
\end{wrapfigure}

詳しく述べよう。物理学にも似たような双対性があり、これは何らかの深い根本原理を指し示しているか、
あるいは私たちの心の働き方について何かを教えてくれている。リチャード・ファインマンは、
量子電磁力学に関する自身の研究においてこの双対性を着想の源として挙げている。

ほとんどの物理法則には二通りの表現形式がある。一つは局所的、あるいは無限小の考察を用いる。
小さな近傍における系の状態を見て、次の瞬間にそれがどのように変化するかを予測する。
これは通常、ある時間にわたって積分、つまり合計されなければならない微分方程式を使って表現される。

このアプローチが命令的な思考に似ていることに気づくだろう。それぞれが前のステップの結果に
依存する一連の小さなステップを辿って最終的な解に到達する。実際、物理系のコンピュータシミュレーションは、
微分方程式を差分方程式に変換して反復するという方法で実装されるのが一般的だ。
アステロイドゲームで宇宙船がアニメーションされる仕組みがこれだ。各タイムステップで、
宇宙船の位置は速度に時間差分を掛けて計算された小さな増分を加えることで変化する。
速度もまた、力を質量で割った加速度に比例する小さな増分で変化する。

これらはニュートンの運動法則に対応する微分方程式を直接コードに落とし込んだものだ：
\begin{align*}
  F & = m \frac{dv}{dt} \\
  v & = \frac{dx}{dt}
\end{align*}
同様の手法は、マクスウェル方程式による電磁場の伝播や、格子\acronym{QCD}
（量子色力学）を用いた陽子内部のクォークとグルーオンの振る舞いといった
より複雑な問題にも適用できる。

デジタルコンピュータの使用によって促進されたこの局所的な考え方と空間・時間の離散化は、
スティーブン・ウルフラムによる宇宙全体の複雑さをセルオートマトンの系に還元しようとする
壮大な試みの中で極限的な表現を得た。

もう一つのアプローチはグローバルなものだ。系の初期状態と最終状態を見て、
ある汎関数を最小化することによってそれらをつなぐ軌跡を計算する。
最も単純な例はフェルマーの最小時間の原理だ。光線は飛行時間を最小化する経路に沿って伝播すると述べている。
特に、反射や屈折を起こす物体がない場合、点$A$から点$B$への光線は最短経路、
つまり直線を取る。しかし光は水やガラスのような密な（透明な）物質中では遅く伝播する。
だから出発点を空気中に、終点を水中に選ぶと、光にとっては空気中を長く進み、
水中で近道をする方が有利になる。最小時間の経路は、光線を空気と水の境界で屈折させ、
スネルの屈折法則をもたらす：
\begin{equation*}
  \frac{sin(\theta_1)}{sin(\theta_2)} = \frac{v_1}{v_2}
\end{equation*}
ここで$v_1$は空気中の光速、$v_2$は水中の光速だ。

\begin{figure}[H]
  \centering
  \includegraphics[width=0.3\textwidth]{images/snell.jpg}
\end{figure}

\noindent
古典力学のすべては最小作用の原理から導出できる。作用は、運動エネルギーと位置エネルギーの差である
ラグランジアンを積分することで任意の軌跡について計算できる（注意：差であり、和ではない――
和は全エネルギーになる）。迫撃砲を発射して目標に当てるとき、発射体はまず上昇する。
そこでは重力による位置エネルギーが高く、作用への負の寄与を積み重ねるためにそこで時間を過ごす。
また放物線の頂点ではゆっくりになり、運動エネルギーを最小化する。
その後、位置エネルギーの低い領域を素早く通過するために加速する。

\begin{figure}[H]
  \centering
  \includegraphics[width=0.35\textwidth]{images/mortar.jpg}
\end{figure}

\noindent
ファインマンの最大の貢献は、最小作用の原理が量子力学にも一般化できることを認識したことだ。
そこでも問題は初期状態と最終状態の観点から定式化される。
それらの状態間のファインマン経路積分が遷移確率の計算に使われる。

\begin{figure}[H]
  \centering
  \includegraphics[width=0.35\textwidth]{images/feynman.jpg}
\end{figure}

\noindent
要点は、物理法則を記述する方法には不思議で未解明の双対性があることだ。
物事が順次、少しずつ起こる局所的な描像を使うこともできる。あるいは初期条件と最終条件を宣言し、
その間のすべてがおのずと従うグローバルな描像を使うこともできる。

グローバルなアプローチはプログラミングでも使われる。例えばレイトレーシングの実装がそうだ。
目の位置と光源の位置を宣言し、それらをつなぐ光線の取りうる経路を求める。
各光線の飛行時間を明示的に最小化するわけではないが、スネルの法則と反射の幾何学を同様の効果のために使う。

局所的アプローチとグローバルなアプローチの最大の違いは、空間、そしてより重要なことには時間の扱い方だ。
局所的アプローチは今ここでの即時の満足を重視するが、グローバルなアプローチは長期的で静的な視点を取る。
まるで未来があらかじめ定められていて、私たちはある永遠の宇宙の性質を分析しているだけであるかのように。

これが最もよく示されているのは、ユーザーインタラクションへの関数型リアクティブプログラミング
（\acronym{FRP}）アプローチだ。可変状態を共有する個々のユーザーアクションごとのハンドラを
別々に書く代わりに、\acronym{FRP}は外部イベントを無限リストとして扱い、
それに一連の変換を適用する。概念的には、私たちの将来のすべてのアクションのリストが、
プログラムへの入力データとして利用可能な形でそこにある。プログラムの視点から見ると、
$\pi$の数字のリスト、擬似乱数のリスト、コンピュータのハードウェアを通じて来るマウス位置のリストの
いずれも違いはない。どの場合も、$n$番目の項目を取得するには、まず最初の$n-1$個の項目を通過しなければならない。
時間的なイベントに適用されるとき、この性質を\emph{因果性}と呼ぶ。

では、これは圏論とどのような関係があるのか？圏論はグローバルなアプローチを促進し、
したがって宣言的プログラミングを支持すると論じたい。まず何より、圏論には微積分とは異なり、
距離、近傍、時間の概念が組み込まれていない。あるのは抽象的な対象とそれらの間の抽象的なつながりだけだ。
$A$から$B$へ一連のステップで到達できるなら、一跳びでそこに到達することもできる。
さらに、圏論の主要なツールは普遍的構成であり、これはグローバルなアプローチの典型だ。
例えば圏論的積の定義でそれが働くのを見た。その性質を指定することで定義された――
非常に宣言的なアプローチだ。それは二つの射影を備えた対象であり、最良の対象だ――
他のそのような対象の射影を因子分解するという性質を最適化する。

\begin{figure}[H]
  \centering
  \includegraphics[width=0.35\textwidth]{images/productranking.jpg}
\end{figure}

\noindent
これをフェルマーの最小時間の原理、あるいは最小作用の原理と比較してみよう。

逆に、これをデカルト積の従来の定義と対比させると、それははるかに命令的だ。
一方の集合から要素を一つ選び、もう一方の集合から別の要素を選ぶことで、
積の要素を生成する方法を記述する。それはペアを作るためのレシピだ。
ペアを分解するための別のレシピもある。

Haskellのような関数型言語を含むほぼすべてのプログラミング言語では、積型、余積型、関数型は
普遍的構成によって定義されるのではなく、組み込みになっている。
ただし、圏論的プログラミング言語を作る試みもある
（例えば\urlref{http://web.sfc.keio.ac.jp/~hagino/thesis.pdf}{萩野達也の論文}を参照）。

直接使われるかどうかに関わらず、圏論的定義は既存のプログラミング構成要素を正当化し、
新たな構成要素を生み出す。最も重要なのは、圏論がコンピュータプログラムについて
宣言的なレベルで推論するためのメタ言語を提供することだ。また、問題仕様をコードに落とす前に
それについて推論することも促進する。
